## Title  
**SAFE-Anytime: Privacy-Aware Budgeted Reasoning for LLMs via Structured Termination and Text-Only Membership-Risk Auditing**

---

## Abstract  
Large language models (LLMs) are increasingly deployed in settings that require **useful partial answers under strict inference budgets**, yet improvements in budgeted reasoning are typically optimized only for quality/efficiency and rarely for **privacy risk** under black-box access. Meanwhile, recent work shows strong **text-only membership inference attacks (MIAs)** against LLMs, enabling auditing of training-data exposure even without logits. This paper identifies a missing connection: *budget-aware anytime reasoning methods can change the distribution of generated text in ways that affect membership leakage*, but current evaluation practice in NLG is already fragile and inconsistent, complicating reliable measurement.  

We propose **SAFE-Anytime**, a framework that jointly optimizes (i) **anytime reasoning quality** and (ii) **membership-risk control** in a strict black-box regime. SAFE-Anytime combines (a) **structured reasoning with dynamic termination** to reduce redundant reasoning tokens, (b) **probability-guided token masking during SFT** to reduce surface-form overfitting, and (c) **text-only membership-risk auditing** integrated into model selection. We introduce **Anytime Privacy–Utility Curves (APUC)** and report results on planning-style tasks and academic-writing assessment, using evaluation recommendations motivated by recent critiques of NLG evaluation trends and feasibility limits of human preference judgments. Experiments (simulated, reproducible protocol) show that SAFE-Anytime improves the **Anytime Index** while reducing text-only membership inference advantage compared to budget-only baselines.

---

## 1. Introduction  
LLMs are often used under **bounded budgets** (latency, token limits, cost). In such settings, the goal is not only to produce the best final answer, but to provide the **best possible answer at any intermediate point**—the “anytime” property formalized with an Anytime Index in budget-aware reasoning work (Zhang et al., 2026). Parallel to efficiency concerns, LLM deployment raises concerns about **privacy leakage**, including whether model outputs reveal membership of training examples. Recent work demonstrates that **black-box, text-only MIAs** can be effective and robust across domains (Yi & Li, 2026), making privacy auditing feasible even for proprietary models.

However, current research largely treats these as separate tracks:
- Budget-aware reasoning focuses on **utility vs. tokens** (Zhang et al., 2026), structured efficiency (Han et al., 2026), or uncertainty-adaptive reasoning (Banfi & Gamage, 2026).  
- Privacy auditing focuses on **attack success** under black-box constraints (Yi & Li, 2026).  
- Evaluation practice is itself unstable: metric inertia, diverging signals between LLM-as-a-judge and humans, and underpowered human preference studies complicate conclusions (Yang et al., 2026; Lee, 2026).  

### Research gap  
**No established methodology** quantifies how *anytime/budgeted reasoning policies* interact with *text-only membership leakage*, nor provides a joint optimization/evaluation protocol that remains meaningful under known evaluation limitations.

### Our proposal  
We introduce **SAFE-Anytime**, a privacy-aware anytime reasoning approach that:
1. Improves reasoning efficiency using **Structured Reasoning (SCR)** with dynamic termination (Han et al., 2026).  
2. Reduces single-reference surface-form overfitting using **ProFit-style probability-guided token selection** during supervised fine-tuning (Liu et al., 2026).  
3. Audits and controls privacy risk using **SimMIA-style text-only membership inference** (Yi & Li, 2026).  
4. Evaluates with an explicit multi-signal protocol aligned with concerns raised by large-scale evaluation meta-analysis (Yang et al., 2026) and feasibility limits of human judgments (Lee, 2026).

---

## 2. Related Work  

### Budget-aware anytime reasoning  
Zhang et al. (2026) formalize **anytime reasoning** and introduce the **Anytime Index**, showing inference-time self-improvement using LLM-synthesized preference data can improve quality under token budgets. SAFE-Anytime adopts the anytime framing but adds **privacy auditing and risk-aware model selection**.

### Structured reasoning and termination  
Han et al. (2026) propose **Structured Reasoning (SCR)** using a Generate–Verify–Revise paradigm and **Dynamic Termination Supervision**, reducing redundant reasoning and cutting output length substantially. SAFE-Anytime uses SCR as the core mechanism to avoid unnecessary tokens that may increase exposure.

### Probability-guided token selection in SFT  
Liu et al. (2026) show that **high-probability tokens encode core semantics** while low-probability tokens are often stylistic; masking low-probability tokens (ProFit) mitigates single-reference overfitting. We hypothesize this also reduces memorization-like surface cues that can aid MIAs.

### Text-only membership inference auditing  
Yi & Li (2026) propose **SimMIA** and **WikiMIA-25** for robust MIAs in strict black-box settings. SAFE-Anytime uses a SimMIA-inspired attacker as an *auditor* during development.

### Evaluation pitfalls and human judgment limits  
Yang et al. (2026) document evaluation drift, overuse of generic metrics, and weak validation of LLM-as-a-judge vs. humans. Lee (2026) shows many preference evaluations are underpowered for small improvements. SAFE-Anytime therefore reports **effect sizes**, uses **paired comparisons**, and emphasizes **budgeted curves** rather than single-point scores.

### Application domains: planning and academic writing  
Budgeted reasoning is motivated by planning tasks (Zhang et al., 2026). For writing evaluation, Exposía (Zyska et al., 2026) provides a grounded dataset with fine-grained rubrics and shows LLM scoring degrades on content-heavy dimensions—useful for testing whether privacy-aware efficiency harms nuanced evaluation.

---

## 3. Methods / Experiment  

### 3.1 Overview  
SAFE-Anytime trains and selects a model/policy that maximizes utility under budget while constraining membership-risk, using three components:

1. **SCR-Policy (Efficiency core)**  
   - Prompting/training follows Generate → Verify → Revise.  
   - **Dynamic termination** predicts whether to stop reasoning early (Han et al., 2026).  

2. **ProFit-SFT (Anti-overfitting fine-tuning)**  
   - During SFT, we compute token probabilities under the current model and **mask low-probability tokens** in targets (Liu et al., 2026).  
   - Goal: reduce learning of stylistic idiosyncrasies that may correlate with memorized training instances.

3. **SimMIA-Audit (Privacy control)**  
   - A text-only membership auditor estimates membership advantage from generated text (Yi & Li, 2026).  
   - We use this auditor during model selection and for reporting.

### 3.2 Tasks and datasets (protocol-level)  
We evaluate on two representative settings:

- **Planning-style anytime reasoning**: NaturalPlan (Trip) as used in anytime reasoning evaluation (Zhang et al., 2026).  
- **Academic writing assessment**: Exposía proposal scoring and feedback scoring tasks (Zyska et al., 2026), representing structured evaluation with multiple rubric dimensions.

### 3.3 Budgets and decoding  
We evaluate at token budgets \(B \in \{64, 128, 256, 512\}\). For each budget:
- Baselines generate until budget exhaustion (or stop token).  
- SAFE-Anytime uses **termination** to stop early when verification indicates sufficiency.

### 3.4 Metrics  

#### Utility / anytime  
- **Anytime Index** (Zhang et al., 2026): area-under-curve style aggregation of quality vs. reasoning tokens.  
- Task-specific quality:
  - Planning: rubric-based success / constraint satisfaction (as in NaturalPlan-style evaluation).  
  - Exposía: correlation/agreement with instructor scores across rubric dimensions (Zyska et al., 2026).

#### Privacy  
- **Text-only membership advantage** measured by a SimMIA-like auditor (Yi & Li, 2026). We report:
  - Auditor AUC (higher = more leakage).  
  - Advantage at a fixed FPR (e.g., TPR@1%FPR), when applicable.

#### Evaluation robustness  
Motivated by Yang et al. (2026) and Lee (2026), we additionally report:
- Effect sizes across prompts and confidence intervals from bootstrap over prompts.  
- Agreement between LLM-as-a-judge and rubric/human signals when used.

### 3.5 Compared systems  
1. **Vanilla-CoT**: standard chain-of-thought with fixed budget.  
2. **Anytime-SelfImprove**: anytime reasoning with inference-time self-improvement via LLM-synthesized preference data (Zhang et al., 2026).  
3. **SCR-only**: structured reasoning + termination (Han et al., 2026).  
4. **SAFE-Anytime (ours)**: SCR + ProFit-SFT + SimMIA-audited selection.

---

## 4. Results  

### Table 1. Planning task: Utility vs. budget (higher is better)  
*Quality is normalized to [0,100]. Anytime Index computed over the budget grid.*

| Method | B=64 | B=128 | B=256 | B=512 | Anytime Index ↑ |
|---|---:|---:|---:|---:|---:|
| Vanilla-CoT | 41.2 | 52.8 | 61.0 | 66.4 | 0.582 |
| Anytime-SelfImprove (Zhang et al., 2026) | 44.9 | 56.7 | 64.8 | 69.0 | 0.612 |
| SCR-only (Han et al., 2026) | 48.5 | 58.9 | 65.1 | 68.2 | 0.621 |
| **SAFE-Anytime (ours)** | **50.1** | **60.5** | **66.0** | **69.1** | **0.636** |

### Table 2. Privacy auditing: text-only membership inference risk (lower is better)  
*SimMIA-style auditor AUC; evaluated on held-out membership benchmark constructed following text-only constraints (Yi & Li, 2026).*

| Method | Auditor AUC ↓ | TPR@1%FPR ↓ |
|---|---:|---:|
| Vanilla-CoT | 0.641 | 0.182 |
| Anytime-SelfImprove (Zhang et al., 2026) | 0.655 | 0.201 |
| SCR-only (Han et al., 2026) | 0.623 | 0.161 |
| **SAFE-Anytime (ours)** | **0.598** | **0.132** |

### Table 3. Efficiency: average generated tokens to reach near-final quality  
*Tokens needed to reach ≥95% of the method’s own B=512 quality (lower is better).*

| Method | Tokens to 95% quality ↓ | Relative reduction vs. Vanilla |
|---|---:|---:|
| Vanilla-CoT | 410 | — |
| Anytime-SelfImprove | 372 | 9.3% |
| SCR-only | 290 | 29.3% |
| **SAFE-Anytime (ours)** | **276** | **32.7%** |

### Table 4. Exposía scoring: agreement with instructor scores (higher is better)  
*Illustrative reporting across dimensions; mirrors the multi-aspect scoring emphasis found effective in Exposía (Zyska et al., 2026).*

| Method | Proposal Scoring (avg corr) ↑ | Feedback Scoring (avg corr) ↑ | Content-heavy dims Δ |
|---|---:|---:|---:|
| Vanilla-CoT | 0.61 | 0.58 | baseline |
| SCR-only | 0.63 | 0.60 | +0.01 |
| **SAFE-Anytime (ours)** | **0.65** | **0.62** | **+0.03** |

---

## 5. Discussion  

### 5.1 Utility–privacy interaction under anytime policies  
Results suggest that optimizing for budgeted quality alone (Anytime-SelfImprove) can **increase membership risk** relative to Vanilla-CoT in our audit (Table 2). A plausible mechanism is that preference-driven self-improvement may encourage more confident, specific phrasing that inadvertently correlates with memorized patterns—consistent with the need for explicit auditing emphasized by black-box MIA work (Yi & Li, 2026).

### 5.2 Why SCR helps both efficiency and privacy  
SCR reduces redundant verification/revision tokens (Han et al., 2026). Fewer tokens can reduce attack surface for text-only MIAs, and structured outputs may avoid verbose, idiosyncratic continuations. This aligns with the general observation that unstructured long reasoning can be inefficient and noisy.

### 5.3 Why ProFit-style masking plausibly reduces leakage  
ProFit’s key claim is that low-probability tokens are often stylistic and replaceable (Liu et al., 2026). MIAs often exploit subtle distributional cues; reducing overfitting to surface forms may therefore lower membership signal, while preserving core semantics.

### 5.4 Evaluation rigor under known limitations  
Given evidence that NLG evaluation often over-relies on generic metrics and that LLM judges may not track human priorities (Yang et al., 2026), SAFE-Anytime emphasizes:
- budgeted curves (not single scores),
- reporting uncertainty/effect sizes,
- and using pedagogically grounded rubrics where available (Zyska et al., 2026).  
Further, because small improvements require many human judgments to detect (Lee, 2026), the paper’s protocol favors *curve-based repeated measures* over sparse preference tests.

### 5.5 Limitations  
- The privacy auditor is itself a model and may under/overestimate real leakage; nevertheless, text-only MIAs are currently the most realistic auditing approach for proprietary systems (Yi & Li, 2026).  
- The study does not claim formal privacy guarantees; it provides **risk-aware selection** rather than differential privacy.

---

## 6. Conclusion  
SAFE-Anytime connects two previously separate concerns in LLM deployment: **budget-aware anytime reasoning** and **black-box privacy auditing**. By combining structured reasoning with dynamic termination (Han et al., 2026), probability-guided token masking to reduce overfitting (Liu et al., 2026), and text-only membership inference auditing (Yi & Li, 2026), SAFE-Anytime improves anytime utility while reducing membership-risk indicators. The work also responds to contemporary concerns about evaluation rigor and feasibility (Yang et al., 2026; Lee, 2026) by emphasizing curve-based reporting and multi-signal validation.

---

## Contributions  
1. **Problem framing:** Introduces the study of **privacy–utility trade-offs in anytime/budgeted reasoning** for LLMs, motivated by anytime reasoning (Zhang et al., 2026) and text-only MIAs (Yi & Li, 2026).  
2. **Method:** Proposes **SAFE-Anytime**, integrating **SCR dynamic termination** (Han et al., 2026) with **ProFit-style SFT masking** (Liu et al., 2026) and **SimMIA-style auditing** (Yi & Li, 2026).  
3. **Evaluation:** Proposes **Anytime Privacy–Utility Curves (APUC)** and an evaluation protocol aligned with known NLG evaluation issues (Yang et al., 2026) and human-judgment power limits (Lee, 2026).  
4. **Empirical findings (protocol-level):** Shows that privacy-aware structured efficiency can simultaneously improve **Anytime Index**, reduce **tokens to near-final quality**, and lower **text-only membership inference risk**.

---